%\documentclass[11pt,a4paper,twoside]{article}

\usepackage[utf8]{inputenc}

\usepackage[T1,T2A]{fontenc}

\usepackage[english.ancient,russian.ancient]{babel}

\usepackage{graphicx}

\usepackage{array}

\usepackage[doi]{samgtu-bib}

\usepackage{samgtu}

\usepackage{samgtu-name}

\usepackage{mathtools}

\mathtoolsset{showonlyrefs}

\usepackage{desclist}

\usepackage{euscript}

\usepackage{mathrsfs} 

\usepackage{multicol}

\usepackage{mathrsfs}

\usepackage{cite}

\usepackage{cmap}

\usepackage{qrcode}

\allowdisplaybreaks[4]

\usepackage[nodayofweek]{datetime}

\usepackage[thinlines]{easybmat}



\renewcommand{\JournalYear}{2018}

\renewcommand{\JournalVolume}{22}

\renewcommand{\JournalNumber}{4}



\newdate{JournalDateSign}{29}{12}{2018}% Дата подписания Вестника в печать

\renewcommand{\JournalDateSign}{\displaydate{JournalDateSign}}

\newdate{JournalDatePrint}{16}{01}{2019}% Дата выхода Вестника из печати

\renewcommand{\JournalDatePrint}{\displaydate{JournalDatePrint}}

%\renewcommand{\JournalUPL}{16.56} % Усл. печ. л.

%\renewcommand{\JournalUIL}{16.53} % Уч.-изд. л.

\renewcommand{\JournalUPL}{15.24} % Усл. печ. л.

\renewcommand{\JournalUIL}{15.21} % Уч.-изд. л.

\renewcommand{\JournalRegNumber}{220/18} % Регистрационный номер

\renewcommand{\JournalOrderNumber}{2} % Номер заказа



%\renewcommand{\JournalRMonth}{Март} % Месяц выхода Вестника

%\renewcommand{\JournalEMonth}{March} % Месяц выхода Вестника

%\renewcommand{\JournalRMonth}{Июнь} % Месяц выхода Вестника

%\renewcommand{\JournalEMonth}{June} % Месяц выхода Вестника

%\renewcommand{\JournalRMonth}{Сентябрь} % Месяц выхода Вестника

%\renewcommand{\JournalEMonth}{September} % Месяц выхода Вестника

\renewcommand{\JournalRMonth}{Декабрь} % Месяц выхода Вестника

\renewcommand{\JournalEMonth}{December} % Месяц выхода Вестника



\newcommand{\totop}[1]{\raisebox{6pt}[1pt][2pt]{#1}}

\newcommand{\tottop}[1]{\raisebox{12pt}[1pt][2pt]{#1}} 

\newcommand{\totttop}[1]{\raisebox{18pt}[1pt][2pt]{#1}} 

\newcommand{\tottttop}[1]{\raisebox{24pt}[1pt][2pt]{#1}} 

\newcommand{\totttttop}[1]{\raisebox{30pt}[1pt][2pt]{#1}} 

\newcommand{\tobot}[1]{\raisebox{-6pt}[1pt][2pt]{#1}} 

\newcommand{\tobbot}[1]{\raisebox{-12pt}[1pt][2pt]{#1}} 

\newcommand{\tobbbot}[1]{\raisebox{-18pt}[1pt][2pt]{#1}}



\graphicspath{{./00PIC/}}

%

%

%\usepackage{samgtu-name}

%

\vsgtu{1579} %registration number

\FirstPage{760}

\LastPage{772}

%

%\pdfcrop

%\draft

%

\ShortCommunication

%

%\begin{document}

%







\hfill \qrcode[hyperlink,height=.8in]{http://doi.org/10.14498/vsgtu1579}

\vspace{-.9in}





\udc{517.958:539.3(1)}

\msc{74F05, 74K20}



\rutitle{Динамическая устойчивость нагретых геометрически~нерегулярных пластин на~основе~модели~Рейснера}

\rutitleshort{Динамическая устойчивость нагретых геометрически нерегулярных пластин\dots}



\entitle{Dynamic stability of heated geometrically irregular plates on the basis of the Reisner model}





\ruauthor{О.~А.~Мыльцина, А.~В.~Полиенко, Г.~Н.~Белосточный}{М\,ы\,л\,ь\,ц\,и\,н\,а~О.~А.,\ П\,о\,л\,и\,е\,н\,к\,о~А.~В.,\ Б\,е\,л\,о\,с\,т\,о\,ч\,н\,ы\,й~Г.~Н.}

\enauthor{O.~A.~Myltcina, A.~V.~Polienko, G.~N.~Belostochny}{M\,y\,l\,t\,c\,i\,n\,a~O.~A.,\

P\,o\,l\,i\,e\,n\,k\,o~A.~V.,\ B\,e\,l\,o\,s\,t\,o\,c\,h\,n\,y~G.~N.}



\ruaffil{\saratovuni.}



\enaffil{\engsaratovuni.}



\ruauthorsdetails{3}{% Количество авторов

\fullauthor{\rupid{98283}{Ольга Анатольевна Мыльцина}}

\CAuthor % Автор-корреспондент

\orcid{0000-0003-4718-2772} % ORCID ID автора 

\regalia{кандидат физико-математических наук} 

\position{ассистент} 

\dept{каф. теории функций и стохастического анализа} 

\email{omyltcina@yandex.ru}

\fullauthor{\rupid{134918}{Асель Валерьевна Полиенко}}

\orcid{0000-0001-6949-4174} % ORCID ID автора 

\position[n]{ведущий инженер} 

\dept{образовательно-научный институт наноструктур и биосистем} 

%\email{aristambecovaav@mail.ru}

\fullauthor{\rupid{98285}{Григорий Николаевич Белосточный}}

%\CAuthor % Автор-корреспондент

\orcid{0000-0003-4471-6599} % ORCID ID автора 

\regalia{доктор технических наук, профессор} 

\position{профессор} 

\dept{каф. математической теории упругости и биомеханики} 

\email{belostochny@mail.ru}

}



%Olga Myltcina https://orcid.org/0000-0003-4718-2772 

%Grigory Belostochny https://orcid.org/0000-0001-6949-4174

%Acel Polienko https://orcid.org/0000-0003-4471-6599







\enauthorsdetails{3}{% Количество авторов

\fullauthor{\rupid{98283}{Olga A. Myltcina}}

\CAuthor % Автор-корреспондент

\orcid{0000-0003-4718-2772} % ORCID ID автора 

\regalia{Cand. Phys. \& Math. Sci.} 

\position{Assistant} 

\dept{Dept. of Functions \& Approxmation Theory} 

\email[n]{omyltcina@yandex.ru}

\fullauthor{\rupid{134918}{Acel V. Polienko}}

\orcid{0000-0001-6949-4174} % ORCID ID автора 

\position[n]{Senior Engineer} 

\dept{Institute of Nanostructures \& Biosystems} 

%\email[n]{aristambecovaav@mail.ru}

\fullauthor{\rupid{98285}{Grigory N. Belostochny}}

%\CAuthor % Автор-корреспондент

\orcid{0000-0003-4471-6599} % ORCID ID автора 

\regalia{Dr. Techn. Sci.} 

\position{Professor} 

\dept{Dept. of Mathematic Theory of Elasticity \&  Biomechanics} 

\email[n]{belostochny@mail.ru} \vspace{-3mm}

}











\rukeywords{термоустойчивость, динамика, нерегулярность, сингулярный, области неустойчивости, континуальный, модель Рейснера, уравнение Матье}

\enkeywords{thermal stability, dynamics, irregularity, singularity, instability areas, Reissner model, Mathieu equation}



\ruabstract{На~основе континуальной модели геометрически нерегулярной пластинки в рамках модели типа Рейснера решается задача динамической устойчивости нагретой ребристой пластинки под действием периодических по временной координате тангенциальных усилий. Тангенциальные усилия в уравнениях динамической устойчивости нагретой пластины конкретизируются на основании решения неоднородной краевой задачи безмоментной термоупругости в перемещениях. Система сингулярных уравнений динамической устойчивости записана в функции прогиба и дополнительных функциях, характеризующих закон изменения касательных напряжений в вертикальных плоскостях по переменным $x$ и $y$. Решение сводится к уравнению Матье, параметры которого представлены в терминах классической теории пластин и~содержат поправки от температуры, поперечных сдвигов и подкрепляющих ребер. Определяются первые три области динамической устойчивости термоупругой системы. Проводится количественный анализ влияния температуры, деформации сдвига в вертикальных плоскостях и относительной высоты ребер на конфигурацию областей динамической устойчивости.}



\enabstract{On the basis of the continuum model of geometrically irregular plate the problem of dynamic stability has been solved. The Reissner type model is considered. The heated plate with ribs is subjected to periodic temporary coordinate of the tangential forces. For the tangential forces a non-homogeneous boundary problem of membrane thermoelasticity in displacements is solved. The system of singular equations of dynamic stability recorded through the function of the deflection and additional functions. The additional functions characterize the law of change of stresses in vertical planes dependent variables $x$ and $y$. The solution is reduced to the Mathieu equation. The characteristics of the Mathieu equation represented by terms in classical theory of plates and contain corrections of temperature, transverse shear and ribs. Three areas of dynamic stability of the thermoelastic system are determined. Quantitative analysis has been carried out. Dependence of the configuration of the areas of dynamic stability on temperature, shear deformation in vertical planes and relative height of ribs is presented.}



\newdate{mylca}{21}{11}{2017}

\date{\displaydate{mylca}}

\newdate{mylcb}{12}{12}{2017}

\revisiondate{\displaydate{mylcb}}

\newdate{mylcc}{18}{12}{2017}

\accepteddate{\displaydate{mylcc}}



\newdate{mylce}{28}{12}{2017}

\renewcommand{\JournalDataOnline}{\displaydate{mylce}}



%\ForPeerReview

\makerutitle













%%%%%%%%%%%%%%%%%%%%%%%%%%%%%%%%%%%%%%%%%%%%%%%%%%%%%%%%%%%%%%%%%%%%%%%%%%%%%%%%%%%%





Рассмотрим геометрически нерегулярную трансверсально изотропную прямоугольную пластинку со сторонами $a$ и $b$, стандартным образом отнесенную к декартовым координатам с началом в левом нижнем углу пластинки [\citen{mylpolbel:bib:Belostochny2, mylpolbel:bib:Belostochny3, mylpolbel:bib:BeloCvet}].

Термоупругая система <<пластинка-ребра>> нагрета до постоянной температуры $\theta_0$, два противоположных края, расположенных по координатным прямым $x=a$ и $y=b$, находятся под действием периодического по временной координате $t$ силового воздействия в виде нестационарной равномерно распределенной нагрузки заданной интенсивности

\begin{equation}\label{eq:myl:f01}

    p(t)=p_1 \cos \vartheta t,

\end{equation}

где $p_1$, $\vartheta$ "--- известные амплитудно-частотные характеристики воздействия. 

Решение динамической термоустойчивости геометрически нерегулярной пластины на основе модели типа Рейснера [\citen{mylpolbel:bib:Ambsrtcumian, mylpolbel:bib:Reissner}] сводится к интегрированию системы сингулярных дифференциальных уравнений, полученных вариационным путем на базе континуальной модели тонкостенных упругих систем~[\citen{mylpolbel:bib:Gilin1, mylpolbel:bib:belosul, mylpolbel:bib:belosmy, mylpolbel:bib:Belostochny1}]:

\begin{multline}

    \Psi_{1\bigcom 1}+\Psi_{2\bigcom 2}+\sum_{i=1}^n \frac{h_1}{h} a_i \Psi_{2\bigcom 2} \delta(x-x_i)=

-\frac{12}{h^3} T^{22}\biggl(

    1+\sum_{i=1}^n \frac{h_i}{h} a_i \delta(x-x_i)\biggr)w_{\bigcom 22} -

    \\ -\frac{12}{h^3} T^{11}w_{\bigcom 11}    +\frac{12 \gamma}{g h^2}\biggl(1+\sum_{i=1}^n \frac{h_i}{h} a_i\delta(x-x_i)\biggr)w_{\bigcom tt},

\end{multline}

\begin{multline}

\label{eq:myl:f02}

(\nabla^2 w)_{\bigcom 1}-\frac{h^2}{10 G'}

\Bigl(\Psi_{1 \bigcom 11}+\frac{1-\nu}{2} \Psi_{1 \bigcom 22}+\frac{1+\nu}{2}\Psi_{2 \bigcom 12}\Bigr)+

\frac{h^3}{12 D}\Psi_1+ \\

 +2({1-\nu})\sum_{i=1}^n \Phi_{3i}

 \Bigl(\frac{h_i}{h}\Bigr)^3 a_i \Bigl(w_{\bigcom 1}-\frac{h^2}{10 G'}\Psi_1\Bigr)_{\bigcom 22}\delta(x-x_i)=0,

\end{multline}

\begin{multline}

(\nabla^2 w)_{\bigcom 2}-\frac{h^2}{10 G'}

\Bigl(\Psi_{2 \bigcom 22}+\frac{1-\nu}{2} \Psi_{2 \bigcom 11}+\frac{1+\nu}{2}\Psi_{1\bigcom 12}\Bigr)+

\frac{h^3}{12 D}\Psi_2+

\\

+\sum_{i=1}^n \Bigl(\Phi_{3i}\Bigl(\frac{h_i}{h}\Bigr)^3 a_i \Bigl(w_{\bigcom 2}-\frac{h^2}{10 G'}\Psi_2\Bigr)_{\bigcom 22}  \Bigr)\delta(x-x_i)=0,

\end{multline}

где $\delta(x-x_i)$ "--- дельта"=функция Дирака \cite{mylpolbel:bib:Mikucincki};   

$\Psi_j(x,y,t),$ $j=1, 2$ "--- искомые функции, связанные с обобщенными углами поворота нормали к~срединной поверхности пластинки  равенствами [\citen{mylpolbel:bib:Ambsrtcumian, mylpolbel:bib:Rassudov}] $$\gamma_j=w_{\bigcom j}-\frac{h^2}{10 G'}\Psi_j;$$ 

$h$ "--- толщина пластинки; 

$G'$ "--- модуль сдвига в плоскостях, перпендикулярных к срединной плоскости пластинки; 

$w(x,y,t)$ "--- функция прогиба; 

$$\Phi_{3i}=1+3\frac{h}{h_i}+3\Bigl(\frac{h}{h_i}\Bigr)^2;$$ 

$n$ "--- число ребер; $h_i$ и $a_i$ "--- соответственно высота и ширина подкрепляющих ребер; 

$\gamma$ "--- удельный вес; 

$g$ "--- интенсивность поля тяжести; 

$$D=\frac{E h^3}{12(1-\nu^2)};$$ 

$E$ "--- модуль упругости; 

$\nu$ "--- коэффициент Пуассона, индексы 1 и 2 после запятой означают частные производные от функций по переменным \emph{x} и \emph{y} соответственно.





Входящие в динамические уравнения \eqref{eq:myl:f02} усилия $T^{11}$ и $T^{22}$ предварительно определяются на основании решения безмоментной неоднородной краевой задачи геометрически нерегулярной пластинки:

\begin{equation}\label{eq:myl:f03}

\begin{split}

&u_{\bigcom 11}+\frac{1-\nu}{2} u_{\bigcom 22}+\frac{1+\nu}{2}v_{\bigcom 12}+ \frac{1-\nu}{2} \sum_{i=1}^n \frac{h_{i}}{h}a_i(u_{\bigcom 2}+v_{\bigcom 1})_{\bigcom 2} \delta(x-x_i)=0,\\

&v_{\bigcom 22}+\frac{1-\nu}{2}v_{\bigcom 11}+\frac{1+\nu}{2}u_{\bigcom 12}+\sum_{i=1}^{n} \frac{h_i}{h} a_i (v_{\bigcom 2}+\nu u_{\bigcom 1})_{\bigcom 2} \delta(x-x_i)=0.

\end{split}

\end{equation}



Сингулярные уравнения \eqref{eq:myl:f03} записаны на основании стандартного предположения о~неучете тангенциальных сил инерции [\citen{ mylpolbel:bib:Ogibalov1, mylpolbel:bib:Ogibalov2}] и постоянной, как отмечалось выше, температуры. Решение этих уравнений, тождественно удовлетворяющих краевым условиям

\begin{equation*}

\begin{array}{lrl}

\text{при} \,\,x=0, \,  x=a{:} & T^{12}=0, & T^{11}=-p(t),\\

\text{при} \,\, y=0,\,\, y=b{:} & v=0, & T^{12}=0,

\end{array}

\end{equation*}

которые перепишутся в компонентах поля перемещений в виде

\begin{equation}

\label{eq:myl:f04}

\begin{array}{lrl}

\text{при}\,\,x=0, \,  x=a{:} & u_{\bigcom 2}+v_{\bigcom 1}=0, & u_{\bigcom 1}+\nu v_{\bigcom 2}=\alpha(1+\nu) \theta_0-\frac{p(t)}{B},\\

\text{при}\,\,y=0, \,\, y=b{:} & v=0, & u_{\bigcom 2}+v_{\bigcom 1}=0,

\end{array}

\end{equation}

запишутся в элементарных функциях:

\begin{equation}\label{eq:myl:f05}

v=0, \quad  u=\alpha(1+\nu) \theta_0 x-\frac{p(t)}{B}x.

\end{equation}



Тангенциальные усилия в системе \eqref{eq:myl:f02} на основании решений \eqref{eq:myl:f05} примут вид

\begin{equation}\label{eq:myl:f06}

    T^{12}=0, \quad  T^{11}=-p(t), \quad  T^{22}=-(E h \alpha \theta_0+\nu p(t)).

\end{equation}

Здесь $B={E h}/({1-\nu^2})$, $\alpha$ "---  коэффициент линейного расширения материала при нагреве, $p(t)$ определяется формулой \eqref{eq:myl:f01}.



Решения системы дифференциальных уравнений \eqref{eq:myl:f02}, тождественно удовлетворяющие условиям шарнирного закрепления сторон пластинки

\begin{equation}\label{eq:myl:f07}

\begin{array}{lrll}

\text{при}\,\,x=0, \,  x=a{:} & w=0, & M^{11}=0, & \Psi_2=0,\\

\text{при}\,\,y=0, \,\, y=b{:} & w=0, & M^{22}=0, & \Psi_1=0,

\end{array}

\end{equation}

зададим в виде

\begin{equation*}

    w(x,y,t)=\sum_{km} w_{km}(t)\sin\frac{k\pi x}{a}\sin\frac{m \pi y}{b},

\end{equation*}

\begin{equation}\label{eq:myl:f08}

    \Psi_1(x,y,t)=\sum_{km} \Psi_{km}^1(t) \cos\frac{k \pi x}{a} \sin \frac{m \pi y}{b},

\end{equation}

\begin{equation*}

        \Psi_2(x,y,t)=\sum_{km} \Psi_{km}^2(t) \sin\frac{k \pi x}{a} \cos \frac{m \pi y}{b},

\end{equation*}

где $M^{jj},$ $j=1,2$ "--- изгибающие моменты, связанные с обобщенными углами поворота равенствами

\begin{equation*}

    M^{jj}=-D(\gamma_{j\bigcom j}+\nu \gamma_{l\bigcom l}), \quad j\neq l=1, 2.

\end{equation*}



Подстановка \eqref{eq:myl:f08} в уравнения системы \eqref{eq:myl:f02} с использованием процедуры Галеркина по пространственной переменной $x$ приводит к системе трех уравнений относительно коэффициентов $w_{km}(t)$, $\Psi_{km}^j(t)$.

Два последних уравнения этой системы (не содержат производных по временной координате) перепишем в виде

\begin{equation}\label{eq:myl:f09}

\begin{split}

  &a_{11} \Psi_{km}^1 + a_{12} \Psi_{km}^2=b_1 w_{km},\\

   &a_{21} \Psi_{km}^1 + a_{22} \Psi_{km}^2=b_2 w_{km},

\end{split}

\end{equation}

где

\begin{equation*}

a_{11}=\frac{h^3}{12 D}

\biggl(\frac{E}{G'}\biggl(\frac{h^2}{10(1-\nu^2)}\Bigl(\Bigl(\frac{k \pi}{a}\Bigr)^2+\frac{1-\nu}{2}

\Bigl(\frac{m \pi}{b}\Bigr)^2\Bigr)+\frac{6}{5}\Bigl(\frac{m \pi h}{b}\Bigr)^2

\sum_{i=1}^n \beta_i^c\biggr)+1\biggr),

\end{equation*}

\begin{multline}

a_{22}=\frac{h^3}{12 D}

\biggl(\frac{E}{G'}\biggl(\frac{h^2}{10(1-\nu^2)}

\Bigl(\Bigl(\frac{m\pi}{b}\Bigr)^2+\frac{1-\nu}{2}\Bigl(\frac{k\pi}{a}\Bigr)^2 \Bigr)+ 

\\

+\frac{(m\pi)^2}{10(1-\nu^2)}\Bigl(\frac{h}{b}\Bigr)^2\sum_{i=1}^n \beta_i^s\biggr)+1\biggr),

\end{multline}

\begin{equation}\label{eq:myl:f10}

    a_{21}=a_{12}=\frac{h^2}{G'}\frac{1+\nu}{20}\frac{km\pi^2}{ab},

\end{equation}

\begin{gather*}

b_1=\Bigl(\frac{k\pi}{a}\Bigr)^3+\frac{k\pi}{a}\Bigl(\frac{m\pi}{b}\Bigr)^2+2(1-\nu)\frac{k\pi}{a}\Bigl(\frac{m\pi}{b}\Bigr)^2 \sum_{i=1}^n \beta_i^c,

\\

b_2=\Bigl(\frac{m\pi}{b}\Bigr)^3+\frac{m\pi}{b}\Bigl(\frac{k\pi}{a}\Bigr)^2+

\Bigl(\frac{m\pi}{b}\Bigr)^3 \sum_{i=1}^n \beta_i^s.

\end{gather*}

Здесь и далее

\begin{equation*}

\beta_i^s=\Phi_{3i}\Bigl(\frac{h_i}{h}\Bigr)^3\frac{a_i}{a}\sin^2 \frac{k \pi x_i}{a}, \quad  \beta_i^c=\Phi_{3i}\Bigl(\frac{h_i}{h}\Bigr)^3\frac{a_i}{a}\cos^2 \frac{k \pi x_i}{a}.

\end{equation*}

Выражения функций $\Psi_{km}^j(t),$ $j=1, 2$, через $w_{km}(t)$ запишутся так:

\begin{equation}\label{eq:myl:f11}

\Psi_{km}^1(t)=\frac{b_1 a_{22}-b_2 a_{12}}{a_{11}a_{22}-a_{21}^2}w_{km}(t), \quad \Psi_{km}^2(t)=\frac{b_2 a_{11}-b_1a_{12}}{a_{11}a_{22}-a_{21}^2}w_{km}(t).

\end{equation}

На основании \eqref{eq:myl:f11} первое уравнение системы, содержащее производную по временной координате, преобразуется к виду

\begin{equation}\label{eq:myl:f12}

\frac{\gamma h a^2}{g}\frac{d^2w_{km}}{dt^2}+

\Bigl((k\pi)^2T^{11}+\Bigl(\frac{m\pi a}{b}\Bigr)^2T^{22}

+ \frac{E h}{12(1-\nu^2)}\Bigl(\frac{h}{a}\Bigr)^2 L_4^{km}\Bigr)w_{km}=0,

\end{equation}

где

\begin{equation}

\label{eq:myl:f13}

L_4^{km}=

\Bigl((k\pi)^2+\Bigl(\frac{m\pi a}{b}\Bigr)^2\Bigr)^2 \widetilde{p}(G',a_i).

\end{equation}

Здесь

\begin{equation}\label{eq:myl:f16}

\widetilde{p}(G',a_i) = \Bigl( {1+\frac{L_1^{km}(G',a_i)}{\bigl((k\pi)^2+\left({m\pi a}/{b}\right)^2\bigr)^2}} \Bigr) \bigl( {1+L_2^{km}(G',a_i)} \bigr)^{-1},

\end{equation}

\begin{equation*}

    \widetilde{p}(G',a_i)=1\quad \text{ при }  G' \rightarrow \infty, \, a_i=0,

\end{equation*}

а $L_j^{km}(G',a_i)$, $j=1,2$ "--- безразмерные величины, содержащие параметры $G'$ и ${a_i}/{a}$, и такие, что $L_j^{km}(G',a_i)=0$ при $G'\rightarrow \infty$, $a_i=0$:

\begin{multline}

\label{eq:myl:f14}

L_1^{km}(G',a_i)=

\Bigl((k\pi)^4+(k\pi)^2\Bigl(\frac{m\pi a}{b}\Bigr)^2\Bigr) \times 

\\ \times

\biggl(\frac{({h}/{a})^2}{10(1-\nu^2)} 

\Bigl(\Bigl(\frac{m\pi a}{b}\Bigr)^2+\frac{1-\nu}{2}(k\pi)^2\Bigr)+\frac{(m \pi)^2}{5(1-\nu^2)}

\Bigl(\frac{h}{b}\Bigr)^2\sum_{i=1}^n \beta_i^s\biggr)

\frac{E}{G'}+\\

+12(1-\nu^2)(k\pi)^2\Bigl(\frac{m\pi a}{b}\Bigr)^2\sum_{i=1}^n\beta_i^c

\biggl(\frac{E}{G'}

\Bigl(\frac{({h}/{a})^2}{10(1-\nu^2)} 

\Bigl(\Bigl(\frac{m\pi a}{b}\Bigr)^2+

\frac{1-\nu}{2}(k\pi)^2\Bigr)+ 

\\

\hspace{6cm}

+

\frac{(m\pi)^2}{5(1-\nu^2)}

\Bigl(\frac{h}{b}\Bigr)^2 \sum_{i=1}^n\beta_i^s\biggr)+1\biggr)-\\

-\biggl(

\Bigl(\frac{m\pi a}{b}\Bigr)^3+\frac{m\pi a}{b}(k\pi)^2+

\Bigl(\frac{m\pi a}{b}\Bigr)^3\sum_{i=1}^n\beta_i^s\biggr)

\frac{km\pi^2}{20(1-\nu^2)}\frac{h}{a}\frac{h}{b}\frac{E}{G'}+ \hspace{2cm} \\

+\Bigl(\Bigl(\frac{m\pi a}{b}\Bigr)^4+(k\pi)^2

\Bigl(\frac{m\pi a}{b}\Bigr)^2\Bigr)

\biggl(\frac{({h}/{a})^2}{10(1-\nu^2)} 

\Bigl((k\pi)^2+\frac{1-\nu}{2}\Bigl(\frac{m\pi a}{b}\Bigr)^2\Bigr)+

\\

\hspace{7cm}   +\frac{6}{5}(m\pi)^2\Bigl(\frac{h}{b}\Bigr)^2   \sum_{i=1}^n \beta_i^c\biggr)\frac{E}{G'}+ \\

+2\Bigl(  \frac{m\pi a}{b}\Bigr)^4 \sum_{i=1}^n \beta_i^s

\biggl( \biggl( \frac{({h}/{a})^2}{10(1-\nu^2)} 

\Bigl((k\pi)^2+\frac{1-\nu}{2}\Bigl(\frac{m\pi a}{b}\Bigr)^2\Bigr)+ \hspace{2cm} 

\\

\hspace{7cm}+\frac{6}{5}(m\pi)^2 \Bigl(\frac{h}{b}\Bigr)^2   \sum_{i=1}^n \beta_i^c\biggr)\frac{E}{G'}+1\biggr)-

\\

-\biggl((k\pi)^3+k\pi\Bigl(\frac{m\pi a}{b}\Bigr)^2+

12(1-\nu^2)k\pi\Bigl(\frac{m\pi a}{b}\Bigr)^2 

\sum_{i=1}^n\beta_i^c \biggr)

\frac{km\pi^2}{20(1-\nu^2)}\frac{h}{a}\frac{h}{b}\frac{E}{G'},

\end{multline}





\begin{multline}

\label{eq:myl:f15}

L_2^{km}(G',a_i)=\frac{E}{G'}

\biggl(

\frac{({h}/{a})^2}{10(1-\nu^2)} 

\Bigl((k\pi)^2+\frac{1-\nu}{2}\Bigl(\frac{m\pi a}{b}\Bigr)^2\Bigr)+ \\

+

\frac{6}{5}(m\pi)^2\Bigl(\frac{h}{b}\Bigr)^2\sum_{i=1}^n \beta_i^c

+\frac{({h}/{a})^2}{10(1-\nu^2)}

\Bigl(\Bigl(\frac{m\pi a}{b}\Bigr)^2+\frac{1-\nu}{2}(k\pi)^2\Bigr)+

\\

\hspace{7.5cm}

+

\frac{(m\pi)^2}{5(1-\nu^2)}

\Bigl(\frac{h}{b}\Bigr)^2\sum_{i=1}^n \beta_i^s\biggr)+

\\

+\Bigl(\frac{E}{G'}\Bigr)^2

\biggl(\Bigl(\frac{({h}/{a})^2}{10(1-\nu^2)} 

\Bigl((k\pi)^2+\frac{1-\nu}{2}\Bigl(\frac{m\pi a}{b}\Bigr)^2\Bigr)+

\frac{6}{5}(m\pi)^2\Bigl(\frac{h}{b}\Bigr)^2\sum_{i=1}^n\beta_i^c\biggr)\times

\\

\times

\biggl(\frac{({h}/{a})^2}{10(1-\nu^2)}

\Bigl(\Bigl(\frac{m\pi a}{b}\Bigr)^2+\frac{1-\nu}{2}(k\pi)^2\Bigr)+

\frac{(m\pi)^2}{5(1-\nu^2)}\Bigl(\frac{h}{b}\Bigr)^2\sum_{i=1}^n \beta_i^s\biggr)-

\\

 -\frac{(km\pi^2)^2(1+\nu)^2}{400(1-\nu^2)^2} \Bigl( \frac{h}{a} \Bigr)^2 \Bigl( \frac{h}{b} \Bigr)^2\biggr).

\end{multline}





%Дробное выражение в формуле \eqref{eq:myl:f13}) обозначим через $\widetilde{p}(G',a_i)$





Для дальнейшего отметим, что убирая в уравнении \eqref{eq:myl:f12} слагаемые, содержащие вторую производную по временной координате, получим уравнение для определения критических сил (или температур) для ребристой пластинки 

на основе модели типа Рейснера, при достижении которых возможен скачкообразный переход к новой форме равновесия с~числом полуволн~$k$ и~$m$, которое запишется следующим образом:

\begin{equation}\label{eq:myl:f17}

(k\pi)^2T^{11}+\Bigl(\frac{m\pi a}{b}\Bigr)^2T^{22}+\frac{E h}{12(1-\nu^2)}

\Bigl(\frac{h}{a}\Bigr)^2L_4^{km}=0.

\end{equation}

Действительно, при тех же краевых условиях получим два случая.



\smallskip



{\small \sc Случай 1.}

При $\theta_0=0$ и $\vartheta=0$: $T^{11}=-p_1$, $T^{22}=-\nu p_1$; 

относительная критическая нагрузка по модели Рейснера вычисляется по формуле

\begin{equation}\label{eq:myl:f18}

\Bigl(\frac{p_1}{E h}\Bigr)_{\text{cr}}^R=

\frac{({h}/{a})^2}{12(1-\nu^2)}

\frac{\bigl((k\pi)^2+({m\pi a}/{b})^2\bigr)^2}{(k\pi)^2+\nu ({m\pi a}/{b})^2} \widetilde{p}(G',a_i).

\end{equation}

Полагая $G'\rightarrow\infty$, $a_i=0$, получим критическую нагрузку для гладкой пластинки на основе модели Лява. Тогда \eqref{eq:myl:f18} примет вид

\begin{equation}\label{eq:myl:f19}

\Bigl(\frac{p_1}{E h}\Bigr)_{\text{cr}}^R=\widetilde{p}(G',a_i)\Bigl(\frac{p_1}{E h}\Bigr)_{\text{cr}}^L.

\end{equation}



\smallskip



{\small \sc  Случай 2.}

Если $\theta_0\neq0$ и $p_1=0$,

то выражение для относительной критической температуры запишется в виде

\begin{equation}\label{eq:myl:f20}

\bigl(\theta_0 \alpha\bigr)_{\text{cr}}^R=

\frac{({h}/{a})^2}{12(1-\nu^2)}

\frac{\bigl((k\pi)^2+({m\pi a}/{b})^2\bigr)^2}{({m\pi a}/{b})^2}\widetilde{p}(G',a_i).

\end{equation}

При тех же предположениях получим

\begin{equation}\label{eq:myl:f22}

\bigl(\theta_0 \alpha\bigr)_{\text{cr}}^R=\widetilde{p}(G',a_i) (\theta_0 \alpha )_{\text{cr}}^L.

\end{equation}



\smallskip



Убирая в уравнении \eqref{eq:myl:f12}) слагаемые с тангенциальными усилиями, получим выражение для квадрата частоты собственных колебаний пластинки в рамках модели Рейснера:

\begin{equation}\label{eq:myl:f22}

\bigl(\omega_{km}^R\bigr)^2=

\frac{E h}{12(1-\nu^2)}\frac{g}{\gamma h a^2}

\Bigl((k\pi)^2+\Bigl(\frac{m\pi a}{b}\Bigr)^2\Bigr)^2 

\widetilde{p}(G',a_i)

\end{equation}

или

$$

\bigl(\omega_{km}^R\bigr)^2=\widetilde{p}(G',a_i) \bigl(\omega_{km}^L\bigr)^2,

$$  

где $\omega_{km}^L$ "--- частота собственных колебаний гладкой пластинки на базе классической теории пластин.



Изложенное выше и вид $p(t)$ позволяют переписать уравнение \eqref{eq:myl:f12}:

\begin{multline}

\label{eq:myl:f23}

\frac{d^2 w_{km}}{dt^2}+\frac{gE}{\gamma a^2}\frac{({h}/{a})^2}{12(1-\nu^2)}

\Bigl((k\pi)^2+\Bigl(\frac{m\pi a}{b}\Bigr)^2\Bigr)^2 

\widetilde{p}\Bigl(\frac{E}{G'}, \frac{a_i}{a}\Bigr)\times\\

 \times\Bigl(1-\frac{\theta_0}{(\theta_0)_{\text{cr}}^R}\Bigr)

 \Bigl(1-\frac{p_1}{(p_1)_{\text{cr}}^R (1- {\theta_0}/{(\theta_0)_{\text{cr}}^R})}\cos\vartheta t\Bigr)w_{km}(t)=0

\end{multline}

или, в стандартных обозначениях, уравнение Матье \eqref{eq:myl:f23} запишется так [\citen{mylpolbel:bib:Ambsrtcumian, mylpolbel:bib:Bolotin}]:

\begin{equation}\label{eq:myl:f24}

    \frac{d^2 w_{km}}{d t^2}+\Omega_{km}^2\left(1-2\mu_{km}\cos\vartheta t\right)w_{km}(t)=0,

\end{equation}

где

\begin{equation}\label{eq:myl:f25}

    \Omega_{km}=\omega_{km}^R\sqrt{1-{\theta_0}/{(\theta_0)_{\text{cr}}^R}} 

\end{equation}

"--- частота собственных колебаний геометрически нерегулярной нагретой пластинки с учетом влияния поперечных сдвигов,

\begin{equation}\label{eq:myl:f26}

    \mu_{km}=\frac{p_1}{2(p_1)_{\text{cr}}^R 

    (1-{\theta_0}/{(\theta_0)_{\text{cr}}^R})}

\end{equation}

"--- коэффициент возбуждения геометрически нерегулярной пластинки с учетом влияния поперечных сдвигов и температуры.



Таким образом, решение динамической задачи нагретой геометрически нерегулярной пластинки под действием периодических по временной координате

силовых воздействий сведено к уравнению Матье [\citen{mylpolbel:bib:Bolotin, mylpolbel:bib:Strett}], которое при некоторых соотношениях между его параметрами $\Omega_{km}$, $\mu_{km}$  и $\vartheta$ допускает неограниченно возрастающие решения, заполняющие целые области~\cite{mylpolbel:bib:Bolotin}.



Определение границ, ограничивающих области динамической неустойчивости с помощью подстановок  для $w_{km}(t)$ вида 

\begin{equation*}

    \sum_{l=1,3,5,\dots}^{\infty}

    \Bigl(a_l \sin\frac{l \vartheta t}{2}+b_l \cos\frac{l \vartheta t}{2}\Bigr) 

\end{equation*}

в случае решения с периодом $2 T$ и

\begin{equation*}

b_0+\sum_{l=2,4,6,\dots}^{\infty}

\Bigl(

a_l \sin\frac{l \vartheta t}{2}+b_l \cos\frac{l \vartheta t}{2}\Bigr) 

\end{equation*}

в случае решения с периодом $T$ \cite{mylpolbel:bib:Bolotin}  в уравнение Матье \eqref{eq:myl:f23},  

позволит с любой точностью определять границы областей динамической неустойчивости. 

Из равенств нулю определителей однородных алгебраических систем для коэффициентов $a_l$, $b_l$ и $b_0$ следуют формулы для границ $\vartheta^*$ первых трех областей динамической неустойчивости [\citen{mylpolbel:bib:Bolotin, mylpolbel:bib:BeloCvet}], которые с учетом полученных формул \eqref{eq:myl:f19}--\eqref{eq:myl:f22}, \eqref{eq:myl:f25}, \eqref{eq:myl:f26} преобразуются к следующему виду:



\begin{itemize}

\item[I] область:

\begin{equation}\label{eq:myl:f27}

\vartheta^*=2 \omega_{km}^L\sqrt{\widetilde{p}-{\theta_0}/{(\theta_0^L)_{\text{cr}}}\pm \mu_{km}^*};

\end{equation}



\item[II] область:



\begin{equation}\label{eq:myl:f28}

\begin{split}

& \vartheta^*=\omega_{km}^L\sqrt{\widetilde{p}-\frac{\theta_0}{(\theta_0^L)_{\text{cr}}}+

\frac13\frac{(\mu_{km}^*)^2}{\widetilde{p}-{ \theta_0}/{(\theta_0^L)_{\text{cr}}}}}, \\

& \vartheta^*=\omega_{km}^L\sqrt{\widetilde{p}-\frac{\theta_0}{(\theta_0^L)_{\text{cr}}}-

2 \frac{ (\mu_{km}^*)^2}{\widetilde{p}-{\theta_0}/{(\theta_0^L)_{\text{cr}}}}};

  \end{split}

\end{equation}



\item[III] область:



\begin{equation}\label{eq:myl:f29}

\vartheta^*=\frac{2}{3}\omega_{km}^L\sqrt{\widetilde{p}-\frac{\theta_0}{ (\theta_0^L)_{\text{cr}}}- 

\frac{9(\mu_{km}^*)^2}{8\left(\widetilde{p}-{\theta_0}/{(\theta_0^L)_{\text{cr}}}\right)\pm 9\mu_{km}^*}},

\end{equation}



\end{itemize}



\noindent 

где $\mu_{km}^*={p_1}/({2(p_{1}^L)_{\text{cr}}})$ "--- коэффициент возбуждения для гладкой ненагретой пластинки ($\theta_0=0$) на основе классической теории гладких пластин.







\begin{figure}[p] \centering



\includegraphics[width=\textwidth]{ris_mylpolbel}



\caption{\label{myl:pic1} Области динамической неустойчивости (залиты цветом; см. формулы (\ref{eq:myl:f27})--(\ref{eq:myl:f29})) квадратной пластины (${a/b=1}$) при различных параметрах; расчеты выполнены для $x_1\hm={a}/{2}$, ${a_1}/{a}=0.01$, ${a}/{h}=50$, $\nu=0.3$  

[Figure~\ref{myl:pic1}. Dynamic instability regions (flooded with color; see Eqs. (\ref{eq:myl:f27})--(\ref{eq:myl:f29})) of a square plate (${a/b=1}$) at various parameters; the calculations are performed for the following values of parameters: $x_1={a}/{2}$, ${a_1}/{a}=0.01$, ${a}/{h}=50$, $\nu=0.3$] }



%  % Requires \usepackage{graphicx}

%\parbox[c]{0.49\linewidth}{

%\centering

% \includegraphics[scale=0.7]{Ris_1_mylpolbel.eps} \\

% \caption{$n=0$, $\frac{E}{G'}=50$, $\frac{\theta_0}{(\theta_0^L)_{\text{кр}}}=0,02$ $\quad$ $\quad$ [Figure 1. $n=0$, $\frac{E}{G'}=50$, $\frac{\theta_0}{(\theta_0^L)_{\text{кр}}}=0,02$]}\label{mylpolbel01}

%}

%\parbox[c]{0.49\linewidth}{

%\centering

% \includegraphics[scale=0.7]{Ris_2_mylpolbel.eps} \\

% \caption{$n=1$, $E/G'=50$, $\frac{\theta_0}{(\theta_0^L)_{\text{кр}}}=0,02$ [Figure 2. $n=1$, $E/G'=50$, $\frac{\theta_0}{(\theta_0^L)_{\text{кр}}}=0,02$]}\label{mylpolbel02}

%}

%\end{figure}

%

%\begin{figure}[h] \centering

%  % Requires \usepackage{graphicx}

%\parbox[c]{0.49\linewidth}{

%\centering

% \includegraphics[scale=0.7]{Ris_3_mylpolbel.eps} \\

% \caption{$n=1$, $\frac{E}{G'}=0$, $\frac{\theta_0}{(\theta_0^L)_{\text{кр}}}=0,02$ $\quad$ [Figure 3. $n=1$, $\frac{E}{G'}=0$, $\frac{\theta_0}{(\theta_0^L)_{\text{кр}}}=0,02$]}\label{mylpolbel03}

%}

%\parbox[c]{0.49\linewidth}{

%\centering

% \includegraphics[scale=0.7]{Ris_4_mylpolbel.eps} \\ \caption{$n=1$, $\frac{E}{G'}=2(1+\nu)$, $\frac{\theta_0}{(\theta_0^L)_{\text{кр}}}=0,02$ $\quad$ [Figure 4. $n=1$, $\frac{E}{G'}=2(1+\nu)$, $\frac{\theta_0}{(\theta_0^L)_{\text{кр}}}=0,02$]}\label{mylpolbel04}

%}

%\end{figure}

%

%\begin{figure}[h] \centering

%  % Requires \usepackage{graphicx}

%\parbox[c]{0.49\linewidth}{

%\centering

% \includegraphics[scale=0.7]{Ris_5_mylpolbel.eps} \\ \caption{$n=0$, $\frac{E}{G'}=50$, $\frac{\theta_0}{(\theta_0^L)_{\text{кр}}}=0$ $\quad$ $\quad$ [Figure 5. $n=0$, $\frac{E}{G'}=50$, $\frac{\theta_0}{(\theta_0^L)_{\text{кр}}}=0$]}\label{mylpolbel05}

%}

%\parbox[c]{0.49\linewidth}{

%\centering

% \includegraphics[scale=0.7]{Ris_6_mylpolbel.eps} \\ \caption{$n=0$, $\frac{E}{G'}=50$, $\frac{\theta_0}{(\theta_0^L)_{\text{кр}}}=0,2$ $\quad$ $\quad$ [Figure 6. $n=0$, $\frac{E}{G'}=50$, $\frac{\theta_0}{(\theta_0^L)_{\text{кр}}}=0,2$]}\label{mylpolbel06}

%}

\end{figure}



\begin{figure}[t!] \centering



\includegraphics[width=\textwidth]{ris_mylpolbel1}



\caption{\label{myl:pic2} Скелетные линии при различных параметрах; расчеты выполнены для $n=5$, $ x_{i+1}-x_i=a/6$, $\theta_0/(\theta_0^L)_{\rm cr}=0.02$

[Figure~\ref{myl:pic2}. Skeletal lines at various parameters; the calculations are performed for the following values of parameters: $n=5$, $ x_{i+1}-x_i=a/6$, $\theta_0/(\theta_0^L)_{\rm cr}=0.02$] }





\end{figure}



Особенность формул \eqref{eq:myl:f27}--\eqref{eq:myl:f29} заключается в том, что они записаны в~терминах	 классической теории пластин с~учетом поправок на влияние температуры, относительной высоты ребер и~учета поперечных сдвигов при различных соотношениях сторон (${a}/{b}$) и форм прогиба пластины ($k, m$).



Количественные результаты, полученные на основании формул \eqref{eq:myl:f27}--\eqref{eq:myl:f29}, представлены на рис.~\ref{myl:pic1},~\ref{myl:pic2}  в осях $\mu=\mu_{11}^*$ и ${\vartheta}/{\omega}={\vartheta^*}/{\omega_{11}^L}$ для случая квадратной пластины (${a}/{b}=1$), когда по направлениям координатных прямых возникает только по одной полуволне ($k=1$, $m=1$), при следующих значениях параметров: $x_1={a}/{2}$, ${a_1}/{a}=0.01$, ${a}/{h}=50$, $\nu=0.3$. 

Залитые цветом области, ограниченные кривыми \eqref{eq:myl:f27}--\eqref{eq:myl:f29}, "--- области динамической неустойчивости: 

первая область неустойчивости "--- наибольшая верхняя область, соответствующая~\eqref{eq:myl:f27}, 

вторая область расположена ниже и получается по формулам~\eqref{eq:myl:f28}, 

третья область соответствует~\eqref{eq:myl:f29}. 





Из данных, представленных на рисунках, следует, что конфигурации областей динамической неустойчивости малочувствительны к изменению перечисленных параметров. 

Во всех рассмотренных случаях области динамической неустойчивости для подкрепленной пластины располагаются выше областей динамической неустойчивости, полученных для гладкой пластины, при прочих равных условиях. 

С~увеличением параметров $\theta_0/\theta_{\rm cr}$  и $E/G'$  области динамической неустойчивости смещаются вниз по оси ${\vartheta}/{\omega}$ (рис.~\ref{myl:pic1}), что согласуется с результатами, приведенным в работе [\citen{mylpolbel:bib:BelZel}], а также известному факту: величины критических температур, при достижении которых происходит статическая потеря термоустойчивости  геометрически нерегулярных пластин и оболочек, растет с~увеличением параметра $h_i/h$  \cite{mylpolbel:bib:BelRassudov18}. При определенных значениях параметра $h_i/h$  и числа ребер жесткости возможно вырождение областей динамической неустойчивости в скелетные линии~(рис.~\ref{myl:pic2}).













%что области динамической устойчивости сужаются с~увеличением температуры пластинки и~параметра ${E}/{G'}$. 

%Включение ребра ведет к~увеличению области неустойчивости (на графиках это заштрихованные области и область правее вертикальной линии), при прочих равных условиях. 

%Этот факт объясняется тем, что постановка ребер жесткости в тонкостенные конструкции в~виде пластин и оболочек, условия эксплуатации которых предусматривают температурные воздействия, всегда увеличивают прогиб [\citen{mylpolbel:bib:Rassudov, mylpolbel:bib:BelRassudov, mylpolbel:bib:BelZel}].







\AdditionalContent{Конкурирующие интересы}{Заявляем, что в отношении авторства и публикации

этой статьи конфликта интересов не имеем.} 

\AdditionalContent{Авторский вклад и ответственность}{Все  авторы  принимали  участие  в  разработке  концепции  \mbox{статьи}  и  в  написании  рукописи. Авторы несут  полную  ответственность  за  предоставление  окончательной рукописи в печать. Окончательная  версия рукописи была одобрена всеми авторами.} 

\AdditionalContent{Финансирование}{Результаты получены в рамках выполнения государственного задания Минобрнауки России № 9.8570.2017/8.9.}





%

%\newpage



\begin{thebibliography}{99}



\RBibitem{mylpolbel:bib:Belostochny2}

\by Белосточный~Г.~Н.

\paper Об одном варианте модели Рейснера геометрически нерегулярных оболочек с термочувствительной толщиной

\inbook Динамические и технологические проблемы механики конструкций и сплошных сред

\bookinfo Материалы III международного симпозиума

\publaddr М.

\publ ЛАТМЭС МГАТУ

\yr 1997

\pages 19--21



\RBibitem{mylpolbel:bib:Belostochny3}

\by Белосточный~Г.~Н., Шкабров~И.~В.

\paper Основные уравнения несвязной термоупругости оболочки с термочувствительной толщиной

\inbook Динамические и технологические проблемы механики конструкций и сплошных сред

\bookinfo Избранные доклады IV международного симпозиума

\publaddr М.

\publ ГРАФРОС

\yr 1998

\pages 65--69





\RBibitem{mylpolbel:bib:BeloCvet}

\by Белосточный~Г.~Н., Цветкова~О.~А.

\paper Геометрически нерегулярные пластинки под действием периодического по времени температурного поля

\inbook Проблемы прочности элементов конструкций под действием нагрузок и рабочих сред

\publaddr Саратов

\publ Сарат. гос. техн. ун-т

\yr 2002

\pages 64--72





\RBibitem{mylpolbel:bib:Ambsrtcumian}

\by Амбарцумян~С.~А.

\book Теория анизотропных пластин

\publaddr М.

\publ Наука

\yr 1967

\totalpages 266



\Bibitem{mylpolbel:bib:Reissner}

\by Reissner~E.

\paper 

 Reflections on the Theory of Elastic Plates

\jour Appl. Mech. Rev.

\vol 38

\issue 11

\pages 1453--1464 

\yr 1985

\crossref{10.1115/1.3143699}







\RBibitem{mylpolbel:bib:Gilin1}

\by Жилин~П.~А.

\paper Линейная теория ребристых оболочек

\jour Изв. АН СССР. МТТ

\yr 1970

\issue 4

\pages 150--162



\RBibitem{mylpolbel:bib:belosul}

\by Белосточный~Г.~Н., Ульянова~О.~И.

\paper Континуальная модель композиции из оболочек вращения с термочувствительной толщиной

\jour Изв. РАН. МТТ

\yr 2011

\issue 2

\pages 32--40

\elib{http://elibrary.ru/item.asp?id=15785027}



\RBibitem{mylpolbel:bib:belosmy}

\by Белосточный~Г.~Н., Мыльцина~О.~А.

\paper Уравнения термоупругости композиций из оболочек вращения

\jour Вестник СГТУ

\yr 2011

\issue 4(59)

\issueinfo Выпуск~1

\pages 56--64



\RBibitem{mylpolbel:bib:Belostochny1}

\by Белосточный~Г.~Н., Пономарев~В.~А.

\preprint Уточненная теория пластин с быстроизменяющимся профилем поперечного сечения и пластин подкрепленных ребрами жесткости сложного очертания (секвенциальный подход)

\preprintinfo  Деп. в~ВИНИТИ 6.06.83 №~3071--83

\yr 1983

\totalpages 15



\Bibitem{mylpolbel:bib:Mikucincki}

\by Antosik~P.,  Mikusinski~J., Sikorski~R.

\book Theory of distributions. The sequential approach

\zmath{0267.46028}

\publaddr Amsterdam

\publ Elsevier Scientific Publishing Company

\totalpages xiv+273 

\yr 1973



\RBibitem{mylpolbel:bib:Rassudov}

\by Рассудов~В.~М., Красюков~В.~П., Панкратов~Н.~Д.

\book Некоторые задачи термоупругости пластинок и~пологих оболочек

\publaddr Саратов

\publ Сарат. ун-т

\yr 1973

\totalpages 154



\RBibitem{mylpolbel:bib:Ogibalov1}

\by Огибалов~П.~М., Грибанов~В.~Ф.

\book Термоустойчивость пластин и оболочек

\publaddr М.

\publ МГУ

\yr 1958

\totalpages 520



\RBibitem{mylpolbel:bib:Ogibalov2}

\by Огибалов~П.~М.

\book Вопросы динамики и устойчивости оболочек

\publaddr М.

\publ МГУ

\yr 1963

\totalpages 417



\RBibitem{mylpolbel:bib:Bolotin}

\by Болотин~В.~В.

\book Динамическая устойчивость упругих систем

\publaddr М.

\publ Гостехиздат

\yr 1956

\totalpages 600



\Bibitem{mylpolbel:bib:Strett}

\by Strutt~M.~J.~O. 

\book Lame, Mathieu and Related Functions in Physics and Technology 

\publ Springer

\publaddr Berlin

\yr 1932



%\RBibitem{mylpolbel:bib:BelRassudov}

%\by Белосточный~Г.~Н., Рассудов~В.~М.

%\paper  Температурные деформации пологих ортотропных оболочек постоянного кручения, подкрепленных ребрами жесткости

%\inbook Механика деформируемых сред

%\publaddr Саратов

%\publ Сарат. ун-т

%\yr 1979

%\pages 83--91



\RBibitem{mylpolbel:bib:BelZel}

\by Белосточный~Г.~Н., Русина~Е.~А.

\paper  Динамическая термоустойчивость трансверсально-изотропных пластин под действием периодических нагрузок

\inbook Современные проблемы нелинейной механики конструкций, взаимодействующих с агрессивными средами

\bookinfo Сб. науч. тр. межвуз. науч. конф. \nofrills

\publaddr Саратов

\yr 2000

\pages 175--180





\RBibitem{mylpolbel:bib:BelRassudov18}

\by Белосточный~Г.~Н., Рассудов~В.~М.

\paper  О потере устойчивости подкрепленных, жестко заделанных по всему контуру ортотропных пластин с учетом влияния поперечных сдвигов

\inbook Механика деформируемых сред

\bookinfo Межвузовск. научн. сборн. \nofrills

\publaddr Саратов

\publ Сарат. ун-т

\yr 1982

\pages 66--69







\end{thebibliography}





\renewcommand{\baselinestretch}{0.82}



\newpage







\makeentitle





\vspace{-5mm}





\AdditionalContent{Competing interests}{We declare that we have no conflicts of interests with the authorship and publication of this article.} 

\AdditionalContent{Authors' contributions and responsibilities}{Each author has participated in the article concept development and in the manuscript writing. The authors are absolutely responsible for submitting the final manuscript in print. Each author has approved the final version of manuscript.} 

\AdditionalContent{Funding}{The results were obtained within the framework of the state task of Russian Ministry of Education and Science (project no. 9.8570.2017/8.9).}





\EngRef







\begin{thebibliography}{99}



\Bibitem{mylpolbel:bib:Belostochny2:en}

\lang In Russian

\by Belostochny~G.~N.

\paper On one version of the Reisner model of geometrically irregular shells with a temperature-sensitive thickness

\inbook Dynamic and technological problems of the mechanics of structures and continuous media

\bookinfo Materials of the 3rd International Symposium

\publaddr Moscow

\yr 1997

\pages 19--21







\Bibitem{mylpolbel:bib:Belostochny3:en}

\lang In Russian

\by Belostochny~G.~N., Shkabrov~I.~V.

\paper Basic equations of uncoupled thermoelasticity of shells of thermosensitive thickness

\inbook Dynamic and technological problems of the mechanics of structures and continuous media

\bookinfo Selected reports of the 4rd International Symposium

\publaddr Moscow

\yr 1998

\pages 65--69









\Bibitem{mylpolbel:bib:BeloCvet:en}

\lang In Russian

\by Belostochny~G.~N., Tsvetkova~O.~A.

\paper Geometrically irregular plates under the action of a time-periodic temperature field

\inbook Problemy prochnosti elementov konstruktsii pod deistviem nagruzok i rabochikh sred \rm [Problems of strength of structural elements under the influence of loads and working media]

\publaddr Saratov

\publ Saratov State Techn. Univ.

\yr 2002

\pages 64--72







\Bibitem{mylpolbel:bib:Ambsrtcumian:en}

\lang In Russian

\by Ambartsumian~S.~A.

\book Teoriia anizotropnykh plastin \rm [Theory of Anisotropic Plates]

\publaddr Moscow

\publ Nauka

\yr 1967

\totalpages 266



\Bibitem{mylpolbel:bib:Reissner}

\by Reissner~E.

\paper 4 Reflections on the Theory of Elastic Plates

\jour Appl. Mech. Rev.

\vol 38

\issue 11

\pages 1453--1464 

\yr 1985

\crossref{10.1115/1.3143699}







\Bibitem{mylpolbel:bib:Gilin1:en}

\lang In Russian

\by Zhilin~P.~A.

\paper Linear theory of ribbed shells 

\jour Izv. AN SSSR. Mekhanika tverdogo tela 

\yr 1970

\issue 4

\pages 150--162



\Bibitem{mylpolbel:bib:belosul:en}

\paper Continuum model for a composition of shells of revolution with thermosensitive thickness

\by Belostochnyi~G.~N., Ul'yanova~O.~I.

\jour Mech. Solids

\yr  2011

\vol 46

\issue 2

\pages 184--191

\crossref{10.3103/S0025654411020051}

\elib{http://elibrary.ru/item.asp?id=16982711}



\Bibitem{mylpolbel:bib:belosmy:en}

\lang In Russian

\by Belostochny~G.~N., Myl'tsina~O.~A.

\paper Equations of thermoelasticity of compositions of shells of revolution

\jour Vestnik SGTU

\yr 2011

\issue 4(59)

\issueinfo Issue~1

\pages 56--64





\Bibitem{mylpolbel:bib:Belostochny1:en}

\lang In Russian

\by Belostochny~G.~N., Ponomarev~V.~A.

\preprint A refined theory of plates with a rapidly changing cross-sectional profile and plates reinforced with stiffening ribs of complex outline. The sequential approach

\preprintinfo  Preprint No. 3071--83; deposited at VINITI  from 6~June 1983

\yr 1983

\totalpages 15





\Bibitem{mylpolbel:bib:Mikucincki:en}

\by Antosik~P.,  Mikusinski~J., Sikorski~R.

\book Theory of distributions. The sequential approach

\zmath{0267.46028}

\publaddr Amsterdam

\publ Elsevier Scientific Publishing Company

\totalpages xiv+273 

\yr 1973



\Bibitem{mylpolbel:bib:Rassudov:en}

\lang In Russian

\by Rassudov~V.~M., Krasiukov~V.~P., Pankratov~N.~D.

\book Nekotorye zadachi termouprugosti plastinok i~pologikh obolochek \rm [Some problems of thermoelasticity of plates and sloping shells]

\publaddr Saratov

\publ Saratov Univ.

\yr 1973

\totalpages 154





\Bibitem{mylpolbel:bib:Ogibalov1:en}

\lang In Russian

\by Ogibalov~P.~M., Gribanov~V.~F.

\book Termoustoichivost' plastin i obolochek \rm [Thermal Stability of Plates and Shells]

\publaddr Moscow

\publ Moscow State Univ.

\yr 1958

\totalpages 520





\Bibitem{mylpolbel:bib:Ogibalov2:en}

\lang In Russian

\by Ogibalov~P.~M.

\book Voprosy dinamiki i ustoichivosti obolochek \rm [Problems of dynamics and stability of shells]

\publaddr Moscow

\publ Moscow State Univ.

\yr 1963

\totalpages 417





\Bibitem{mylpolbel:bib:Bolotin:en}

\lang In Russian

\by Bolotin~V.~V.

\book Dinamicheskaia ustoichivost' uprugikh sistem \rm [Dynamic Stability of Elastic Systems]

\publaddr Moscow

\publ Gostekhizdat

\yr 1956

\totalpages 600





\Bibitem{mylpolbel:bib:Strett:en}

\by Strutt~M.~J.~O. 

\book Lame, Mathieu and Related Functions in Physics and Technology 

\publ Springer

\publaddr Berlin

\yr 1932



%\Bibitem{mylpolbel:bib:BelRassudov:en}

%\by Belostochny~G.~N., Rassudov~V.~M.

%\paper  Temperature deformations of gently sloping orthotropic constant-torsion shells reinforced by stiffening ribs

%\inbook Mekhanika deformiruemykh sred \rm [Mechanics of deformable media]

%\publaddr Saratov

%\publ Saratov Univ.

%\yr 1979

%\pages 83--91









\Bibitem{mylpolbel:bib:BelZel:en}

\lang In Russian

\by Belostochny~G.~N., Rusina~E.~A.

\paper  Dynamic thermal stability of transversal-isotropic plates under the action of periodic loads

\inbook Sovremennye problemy nelineinoi mekhaniki konstruktsii, vzaimodeistvuiushchikh s agressivnymi sredami

\rm [Modern problems of nonlinear mechanics of structures interacting with aggressive media]

\publaddr Saratov

\yr 2000

\pages 175--180





\Bibitem{mylpolbel:bib:BelRassudov18:en}

\lang In Russian

\by Belostochny~G.~N., Rassudov~V.~M.

\paper  On the loss of stability of reinforced rigidly fixed on the entire contour of orthotropic plates taking into account the effect of transverse shifts

\inbook Mekhanika deformiruemykh sred \rm [Mechanics of deformable media]

\publaddr Saratov

\publ Saratov Univ.

\yr 1982

\pages 66--69





















\end{thebibliography}





\renewcommand{\baselinestretch}{0.9}





%\end{document}









